\chapter{AI Vision}
\label{ch:vision}

The AI Vision sensor is a camera with its own processor. It does not send whole
images to the Brain. Instead, it finds areas of the image that match colours we
have described, and reports a list of \emph{objects}: where each one is and how
big it is. This chapter explains how we describe the cone to the sensor, how
the program picks the cone out of the sensor's reports, and the diagnostic tools
we built when the sensor did not work.

\section{How the Sensor Reports Objects}
\label{sec:sensor}

The sensor's image is 320 pixels wide and 240 pixels tall. For each object it
finds, it reports:

\begin{itemize}
  \item \code{id}: which colour (or colour code) matched;
  \item \code{centerX}, \code{centerY}: the centre of the object, in pixels;
  \item \code{width}, \code{height}: the size of the object's bounding box, the
    smallest rectangle that contains it.
\end{itemize}

These numbers are enough for our purposes. \code{centerX} tells us whether the
cone is to the left (below 160) or right (above 160) of the robot, and
\code{width} tells us roughly how close it is, since a closer cone looks wider.

\section{Describing the Cone}
\label{sec:describe}

Our cone is half yellow and half blue. We give the sensor two \emph{colour
descriptions}, one for each half, and one \emph{colour code}, which matches only
when the two colours are found next to each other. The colour code is more
reliable than either colour alone, because many things on a field are yellow or
blue, but very few are yellow directly beside blue.

\lstinputlisting[firstline=151,firstnumber=151,lastline=160,caption={Colour descriptions in \code{src/config.py}, lines 151--160.},label=lst:colours]{\src{config.py}}

Each \code{Colordesc} gives an id, a red--green--blue (RGB) value, and two
tolerances: how far the hue (the shade of colour) and the saturation (how vivid
it is) may differ from that value and still count as a match. The RGB values
were sampled from the AI Vision Utility, a VEX tool that shows the camera's view
on a laptop. The tolerances are wide (hue 40, saturation 1.00), which helps the
sensor recognise the cone in different lighting, but might also match objects
that are not the cone.

These descriptions are passed to \code{AiVision()}, which sends them to the sensor
when the program starts. Line 160 then switches colour detection on. Both of
these details were discovered the hard way, as described in
\sref{sec:ch-silent}.

\section{Finding the Cone}
\label{sec:findcone}

\code{find_cone()} is the one function that every other part of the program uses
to look for the cone. It returns the cone's object, or \code{None} if there is no
cone in view.

\lstinputlisting[firstline=9,firstnumber=9,lastline=26,caption={\code{find_cone()} in \code{src/vision.py}, lines 9--26.},label=lst:findcone]{\src{vision.py}}

It first asks the sensor for every colour code it can see (line 14). It ignores
any that are not the cone code, and any smaller than 6 by 6 pixels, since a tiny
match is more likely to be noise than a cone. If several cones remain, it
returns the largest, because the largest is most likely the nearest.

The minimum size of 6 pixels was originally 20. We lowered it because a cone on
the far side of the field is only about 15 pixels wide, so the robot could not
see distant cones at all (commit \commit{6c3a1b7}).

\section{The Colour-Pair Fallback}
\label{sec:colourpair}

At a distance or in poor lighting, the sensor sometimes recognises both colours
but fails to recognise the code that combines them. When this happens,
\code{find_cone()} falls back to \code{find_color_pair()}, which does the
combination itself:

\lstinputlisting[firstline=36,firstnumber=36,lastline=56,caption={\code{find_color_pair()} in \code{src/vision.py}, lines 36--56.},label=lst:pair]{\src{vision.py}}

It finds the largest yellow object and the largest blue object. It then checks
that they are close together: their centres must be no further apart than the
largest side of either object (lines 49--51). If they are, it builds one box that
surrounds both (lines 52--56).

This merged box matters. An earlier version returned only the yellow half, so
the robot aimed at the middle of the yellow half instead of the middle of the
cone, and judged distance using half the cone's real width. Merging the two
halves into one box fixed both problems (commits \commit{8471b29} and
\commit{f1f4d77}).

The box is returned as a \code{ConeBox} (lines 28--34), a small class with the
same field names as a real sensor object. This means code that uses
\code{find_cone()} does not need to know which method found the cone.

\section{Confirming a Sighting}
\label{sec:confirm}

A single frame can contain a false match, for example a flash of light that
briefly looks yellow. Acting on every single sighting would make the robot
twitch. The program therefore waits for the cone to appear in
\code{VISION_CONFIRM_FRAMES} (3) frames in a row before believing it is there. At
one frame every 50~ms, this adds only 150~ms of delay. \code{cone_detector()}
(lines 58--98 of \code{vision.py}) uses the same idea in reverse: the cone must be
missing for 3 frames before the robot reports it as lost.

\section{Diagnostic Tools}
\label{sec:diagnostics}

When the sensor did not detect anything, there were several possible causes: a
wrong port, an unplugged cable, the colours not being loaded, or the program not
running at all. We built two tools to tell these apart without a laptop.

\subsection{Vision Report}

\code{vision_report()} draws the sensor's status on the Brain screen four times
per second:

\begin{itemize}
  \item a \textbf{tick counter} that goes up each time it redraws, which proves
    the program is running;
  \item \code{installed()}, which is \code{True} if a sensor is plugged into the
    configured port;
  \item the sensor's \textbf{timestamp}, which changes with every new frame, which
    proves the sensor is sending data;
  \item the number of \textbf{codes} and \textbf{colours} found, and the details of
    each object.
\end{itemize}

Each line rules out one cause. For example, if the tick climbs and the timestamp
changes but the counts stay at zero with the cone in view, then the program,
cable and sensor all work, and the problem must be the colour settings. This is
exactly what the screen showed on 22 September (\sref{sec:ch-silent}).

\subsection{Port Scanner}
\label{sec:portscan}

\code{scan_ports()} finds which port the sensor is on, which is useful because
the port in \code{config.py} was originally a guess:

\lstinputlisting[firstline=150,firstnumber=150,lastline=160,caption={Port probing in \code{scan_ports()}, \code{src/vision.py}, lines 150--160.},label=lst:scan]{\src{vision.py}}

It tries to create a sensor on each of the 21 ports. A port counts as a sensor
only if it reports \code{installed()} \emph{and} its timestamp changes over
100~ms; a port with nothing on it will not produce new frames. The
\code{try}/\code{except} block skips any port that causes an error, so one bad
port does not stop the scan.

\section{Summary}

\begin{enumerate}
  \item The sensor reports each matching object's position and size
    (\sref{sec:sensor}).
  \item The cone is described as a yellow--blue colour code, declared in code and
    sent to the sensor at start-up (\sref{sec:describe}).
  \item \code{find_cone()} returns the largest cone, falling back to merging a
    nearby yellow and blue object (\sref{sec:findcone}, \sref{sec:colourpair}).
  \item On-screen diagnostics and a port scanner find faults without a laptop
    (\sref{sec:diagnostics}).
\end{enumerate}

Having described the final program, we now describe how it developed over
time.
